% Source: https://artofproblemsolving.com/community/c3h2179413_phi_functions
% Stable identity: aops-topic-2179413-q01
\textbf{(a)} An integer from $1$ to $p^\alpha$ fails to be relatively prime to $p^\alpha$ exactly when it shares the prime factor $p$, that is, when it is a multiple of $p$. The multiples of $p$ from $1$ to $p^\alpha$ are $p,2p,\dots,p^{\alpha-1}\cdot p$, which is $p^{\alpha-1}$ numbers. So
\[\varphi\left(p^\alpha\right)=p^\alpha-p^{\alpha-1}=p^\alpha\left(1-\dfrac{1}{p}\right).\]

\textbf{(b)} ``Multiplicative'' means $\varphi(mn)=\varphi(m)\varphi(n)$ \emph{when $m$ and $n$ are relatively prime}. The factors of $p^\alpha=p\cdot p\cdots p$ are not relatively prime, so the rule does not apply. For example, $\varphi(4)=2$ (the numbers $1,3$), but $\varphi(2)^2=1$.

The correct computation applies multiplicativity to the coprime prime powers $2^3$, $3^2$, $5^3$ and then uses (a):
\[\varphi(9000)=\varphi\left(2^3\right)\varphi\left(3^2\right)\varphi\left(5^3\right)=(8-4)(9-3)(125-25)=4\cdot6\cdot100=\boxed{2400},\]
which agrees with $9000\left(1-\dfrac{1}{2}\right)\left(1-\dfrac{1}{3}\right)\left(1-\dfrac{1}{5}\right)=2400$.
